% !TeX program = xelatex
% Full English counterpart of dyson_d26.tex. Compile twice with XeLaTeX.
% All references are embedded; no external figures or BibTeX files are required.
\documentclass[a4paper,11pt]{article}
\usepackage[margin=23mm]{geometry}
\usepackage{amsmath,amssymb,amsthm,mathtools}
\usepackage{booktabs,array,enumitem,xcolor}
\usepackage{float}
\usepackage{titlesec,fancyhdr,listings,needspace,etoolbox}
\usepackage[unicode,colorlinks=true,linkcolor=blue!45!black,
  citecolor=blue!45!black,urlcolor=blue!45!black]{hyperref}
\setlength{\parindent}{1.5em}
\setlength{\parskip}{.3em}
\setlength{\emergencystretch}{2em}
\linespread{1.12}
\renewcommand{\abstractname}{Overview}
\allowdisplaybreaks[2]
\titleformat{\section}{\Large\bfseries}{\thesection}{.8em}{}
\titleformat{\subsection}{\large\bfseries}{\thesubsection}{.7em}{}
\titlespacing*{\section}{0pt}{2.2ex plus 1ex}{1ex}
\titlespacing*{\subsection}{0pt}{1.8ex plus .7ex}{.7ex}
\setlist{leftmargin=2em,itemsep=.25em,topsep=.4em}
\pagestyle{fancy}
\fancyhf{}
\renewcommand{\headrulewidth}{0pt}
\fancyfoot[C]{\thepage}
\setlength{\headheight}{15pt}
\theoremstyle{definition}
\newtheorem{proposition}{Proposition}[section]
\newtheorem{lemma}[proposition]{Lemma}
\newtheorem{corollary}[proposition]{Corollary}
\newtheorem{remark}[proposition]{Remark}
\BeforeBeginEnvironment{proposition}{\Needspace{7\baselineskip}}
\BeforeBeginEnvironment{corollary}{\Needspace{6\baselineskip}}
\DeclareMathOperator{\Repart}{Re}
\DeclareMathOperator{\rank}{rank}
\newcommand{\Z}{\mathbb Z}
\newcommand{\Q}{\mathbb Q}
\newcommand{\C}{\mathbb C}
\newcommand{\T}{\mathcal T}
\newcommand{\norm}{\mathrm N}
\newcommand{\iu}{\mathrm i}
\newcommand{\diff}{\mathrm d}
\lstset{language=Python,basicstyle=\small\ttfamily,
  keywordstyle=\color{blue!55!black},commentstyle=\color{black!60},
  breaklines=true,columns=fullflexible,keepspaces=true,
  showstringspaces=false,frame=single,rulecolor=\color{black!20},
  xleftmargin=.6em,xrightmargin=.6em,aboveskip=1em,belowskip=1em}

\hypersetup{pdftitle={Coefficient Formulas for the 26th Power of Dedekind's Eta Function},
 pdfsubject={Atkin product, harmonic polynomials, CM forms, coefficient sums}}

\title{\textbf{Coefficient Formulas for the 26th Power of\\
Dedekind's \(\eta\)-Function:\\
Atkin's Product, Harmonic Polynomials and CM Decomposition}}
\author{}
\date{10 October 2026}
\begin{document}
\maketitle
\vspace{-1.8em}
\begin{abstract}
The 26th power of Dedekind's eta function has both a twelve-factor formula
at prime indices and sum expressions for all coefficients.
Starting with coefficient indexing and a worked example, these notes
rewrite Atkin's product as a difference of degree-twelve harmonics with equal norms.
A parametrized quadric explains factor pairing, while complex multiplication
explains finite signs and cancellation at composite indices.
The final sections connect the formulas with Jacobi specialization
and give methods for computing and checking them.
\end{abstract}
\noindent\textbf{Keywords:} Dedekind eta function; Atkin formula;
harmonic polynomials; complex multiplication; binary quadratic forms;
coefficient vanishing.

\begin{quote}\small
\textbf{About these notes.} This is a personal study compilation based mainly on
the work of Atkin, Serre, Chan--Cooper--Toh and others.
The changes of variables, derivations and examples are intended to aid understanding;
no priority is claimed for the formulas or proofs. Sources are cited where they are used.
\end{quote}

\section{Starting with the coefficient problem}
\subsection{Three formulas, three viewpoints}
Let \(\tau\) lie in the upper half-plane, set \(q=e^{2\pi\iu\tau}\), and define
\begin{equation}\label{eq:eta}
 \eta(\tau)=q^{1/24}\prod_{j\geq1}(1-q^j).
\end{equation}
We want to understand the coefficients in the expansion of \(\eta(\tau)^{26}\).
There are three useful viewpoints.
Atkin's formula gives a product of twelve linear factors at prime indices.
Serre's complex multiplication (CM) decomposition combines four weighted norm series.
The Chan--Cooper--Toh double series gives all coefficients as signed sums
over binary quadratic forms.
Atkin's product is printed in Dyson's Note 14 \cite{Dyson1972};
the latter two expressions are in \cite{Serre1985,ChanCooperToh2006}.

A useful thread to follow is that the twelve factors come from a difference
of two equal-norm harmonics, congruences control the signs,
and composite indices require adding several representations.
Section 2 begins with the product and harmonic difference;
Section 3 explains factor pairing; Section 4 supplies CM and finite signs;
and Section 5 treats all coefficients and their vanishing.
Sections 6 and 7 connect Jacobi specialization and computational checks.
On a first reading, the formulas and examples in Sections 1, 2 and 5
can be read before returning to the more abstract proofs in Sections 3 and 4.

\subsection{Coefficient indexing and conventions for prime representations}
To keep shifted indices distinct from ordinary power-series indices, write
\begin{align}
 E_{26}(q)&=\prod_{j\geq1}(1-q^j)^{26}
          =\sum_{m\geq0}c_{26}(m)q^m,\label{eq:Euler26}\\
 \eta(\tau)^{26}
 &=\sum_{m\geq0}c_{26}(m)q^{m+13/12}
   =\sum_{\substack{n\geq13\\n\equiv1\ (12)}}
       \tau_{26}(n)q^{n/12}.\label{eq:offset}
\end{align}
Thus
\begin{equation}\label{eq:index}
 n=12m+13,\qquad \tau_{26}(12m+13)=c_{26}(m).
\end{equation}
For example, ordinary indices \(m=0,1,2\) correspond to \(n=13,25,37\), and
\[
 E_{26}(q)=1-26q+299q^2-1950q^3+\cdots.
\]
Thus \(\tau_{26}(13)=c_{26}(0)=1\),
\(\tau_{26}(25)=c_{26}(1)=-26\), and
\(\tau_{26}(37)=c_{26}(2)=299\).
We use \(p\) when the shifted index \(n\) is prime.
Atkin's product applies to such primes \(p\equiv1\pmod{12}\);
the index \(n=25\) instead requires the full coefficient sum.
The integers \(a,b,c,d\) will be coordinates in quadratic-form representations;
\(d\) is simply one of those coordinates.

For a prime \(p\equiv1\pmod{12}\), choose \(a,b,c,d>0\) and impose \(3\mid a\)
to fix the order of the two squares. Then there is a unique representation
\begin{equation}\label{eq:representations}
 p=a^2+b^2=c^2+3d^2,\qquad 3\mid a .
\end{equation}
Positivity fixes the signs, and \(3\mid a\) fixes the order of the two squares.
The quadratic forms are norms in two imaginary quadratic fields:
\[
 \norm_{\Q(\iu)/\Q}(a+\iu b)=a^2+b^2,\qquad
 \norm_{\Q(\sqrt{-3})/\Q}(c+d\sqrt{-3})=c^2+3d^2.
\]
They reflect splitting in the Gaussian and Eisenstein quadratic fields.
The condition \(p\equiv1\pmod{12}\) ensures that \(p\) splits in both.

\section{How the twelve factors arise from a harmonic difference}
\subsection{The formula and its domain of validity}
Define
\begin{align}
 P(a,b,c,d)
 &=(a-c)(a+c)(b-c)(b+c)\notag\\
 &\quad{}\times
   \prod_{\epsilon,\delta=\pm1}(2a+\epsilon c+3\delta d)
   \prod_{\epsilon,\delta=\pm1}(2b+\epsilon c+3\delta d),
   \label{eq:P}\\
 C&=2^3\cdot3^4\cdot11^2\cdot13=1019304.\label{eq:C}
\end{align}
With the conventions of Dyson's Note 14, for primes represented as in
\eqref{eq:representations},
\begin{equation}\label{eq:Atkin}
 \tau_{26}(p)=\frac{(-1)^a}{C}P(a,b,c,d).
\end{equation}
The first two pairs in \eqref{eq:P} each supply two factors, and the two
products each supply four, giving twelve in total.
The notation \(\epsilon,\delta=\pm1\) means that all four sign combinations occur.

\paragraph{A worked example: \(p=13\).}
Since \(13=3^2+2^2=1^2+3\cdot2^2\), take \((a,b,c,d)=(3,2,1,2)\).
Compute the factors in three groups:
\[
 \begin{aligned}
 (a-c)(a+c)(b-c)(b+c)&=2\cdot4\cdot1\cdot3=24,\\
 \prod_{\epsilon,\delta=\pm1}(2a+\epsilon c+3\delta d)
 &=13\cdot1\cdot11\cdot(-1)=-143,\\
 \prod_{\epsilon,\delta=\pm1}(2b+\epsilon c+3\delta d)
 &=11\cdot(-1)\cdot9\cdot(-3)=297.
 \end{aligned}
\]
Hence \(P=24\cdot(-143)\cdot297=-C\).
Multiplying by \((-1)^a=-1\) gives \(\tau_{26}(13)=1\).
The example checks both the constant and the sign.
Composite indices involve multiple representations and cancellation, as discussed in Section 5.

\subsection{What the twelve factors encode}
Define
\begin{equation}\label{eq:H}
 H_{\mathrm G}(a,b)=\Repart(a+\iu b)^{12},\qquad
 H_{\mathrm E}(c,d)=\Repart(c+\iu\sqrt3\,d)^{12}.
\end{equation}
Homogeneity of degree twelve means that scaling all coordinates by \(t\)
scales the value by \(t^{12}\). Harmonicity means that the corresponding
Laplacian is zero. Real parts of twelfth powers have both properties.
For the second polynomial the Euclidean coordinates are \((c,\sqrt3\,d)\), so
\[
 (\partial_a^2+\partial_b^2)H_{\mathrm G}=0,\qquad
 (\partial_c^2+\tfrac13\partial_d^2)H_{\mathrm E}=0.
\]

\begin{proposition}\label{prop:harmonic}
On the quadric \(a^2+b^2=c^2+3d^2\),
\begin{equation}\label{eq:harmonicidentity}
 H_{\mathrm G}(a,b)-H_{\mathrm E}(c,d)=-8P(a,b,c,d).
\end{equation}
\end{proposition}
\begin{proof}
Set \(X=a^2,Y=b^2,U=c^2,V=3d^2\), and write
\[
 n=X+Y=U+V,\qquad T=XY,\qquad S=UV.
\]
The binomial expansion yields
\begin{equation}\label{eq:F}
 \Repart(x+\iu y)^{12}
 =\mathcal H_{x^2+y^2}(x^2y^2),\qquad
 \mathcal H_n(T)=n^6-72n^4T+768n^2T^2-2048T^3 .
\end{equation}
Consequently,
\begin{align}
 \mathcal H_n(T)-\mathcal H_n(S)
 =-8(T-S)\bigl[
 9n^4-96n^2(T+S)+256(T^2+TS+S^2)
 \bigr].\label{eq:Hfactor}
\end{align}
On the other hand, \((X-U)(Y-U)=T-S\), and pairing the factors involving
\(a\) and \(b\) in \eqref{eq:P} gives
\begin{align*}
 \prod_{\epsilon,\delta=\pm1}(2a+\epsilon c+3\delta d)
 &=(4X+U-3V)^2-16XU,\\
 \prod_{\epsilon,\delta=\pm1}(2b+\epsilon c+3\delta d)
 &=(4Y+U-3V)^2-16YU.
\end{align*}
Their product is
\[
 9n^4-96n^2(T+S)+256(T^2+TS+S^2).
\]
Together with \eqref{eq:Hfactor}, this proves the identity.
\end{proof}

The prime formula is therefore equivalent to
\begin{equation}\label{eq:AtkinH}
 \tau_{26}(p)=\frac{(-1)^a}{8C}
       \bigl[H_{\mathrm E}(c,d)-H_{\mathrm G}(a,b)\bigr].
\end{equation}
Equation \eqref{eq:AtkinH} compresses the twelve factors into a difference
of two real parts. The two forms serve different purposes:
the product is convenient for computation and for testing whether a factor vanishes,
whereas the harmonic form displays the common norm and angular structure.
It also connects directly with the CM series and quadratic-form sums
below \cite{Serre1985,ChanCooperToh2006}.

Writing \(a+\iu b=\sqrt p\,e^{\iu\theta}\) and
\(c+\iu\sqrt3\,d=\sqrt p\,e^{\iu\varphi}\) gives
\begin{equation}\label{eq:angular}
 \tau_{26}(p)=\frac{(-1)^a p^6}{8C}
       \bigl[\cos(12\varphi)-\cos(12\theta)\bigr].
\end{equation}
The common norm \(p\) supplies the same radial factor \(p^6\);
the two representation angles enter the cosine difference.
Here twelve is first an angular harmonic degree.
The unit groups have orders 4 and 6, so twelfth powers remove both unit phases.
The full coefficient also needs the congruence signs carried by finite characters,
which Section 4 explains.

\Needspace{10\baselineskip}
\subsection{Why these prime coefficients are nonzero}
\begin{corollary}\label{cor:prime}
For every prime \(p\equiv1\pmod{12}\), \(\tau_{26}(p)\ne0\).
\end{corollary}
\begin{proof}
If \(a=\pm c\), equality of norms gives \(b^2=3d^2\), impossible for
positive integers \(b,d\). If \(b=\pm c\), it similarly gives
\(a^2=3d^2\).
If \(2a+\epsilon c+3\delta d=0\), substitution into the norm equation
yields
\[
 b^2=3(a+2\delta d)^2.
\]
Integer solutions would require \(b=0\), again impossible.
The four factors involving \(b\) are handled in the same way.
All factors of \eqref{eq:P} are nonzero, and the assertion follows
from Atkin's formula.
\end{proof}
This argument uses the canonical prime representation and the twelve-factor product.
Composite coefficients may also involve cancellation; Section 5 gives an explicit zero coefficient.


\section{How the factors are paired}
\subsection{Two families of one-variable factors on the quadric}
The identity in the preceding section holds on the equal-norm quadric.
For instance, at \((a,b,c,d)=(1,0,0,0)\),
\(H_{\mathrm G}-H_{\mathrm E}=1\), whereas \(P=0\).
The constraint must therefore be retained when discussing factorization.
Make the invertible complex linear change of coordinates

\begin{equation}\label{eq:complexcoords}
 x=a+\iu b,\quad y=a-\iu b,\quad
 u=c+\iu\sqrt3\,d,\quad v=c-\iu\sqrt3\,d.
\end{equation}
\begin{equation}\label{eq:quadric}
 Q_0=a^2+b^2-c^2-3d^2=xy-uv.
\end{equation}
To find coordinates suited to factorization, arrange the four entries in a matrix.
The equation \(xy-uv=0\) says that its rank is at most one,
so it can be written as a column vector times a row vector:
\[
 \begin{pmatrix}x&u\\v&y\end{pmatrix}
 =\begin{pmatrix}r_0\\r_1\end{pmatrix}
  \begin{pmatrix}s_0&s_1\end{pmatrix}.
\]
Comparing entries gives
\begin{equation}\label{eq:Segre}
 x=r_0s_0,\quad u=r_0s_1,\quad
 v=r_1s_0,\quad y=r_1s_1.
\end{equation}

Equation \eqref{eq:Segre} is the Segre parametrization.
Substituting these expressions into a polynomial is called pulling it back
to the parameter space. Geometrically, this is related to the two families
of lines on the projective quadric; see Sections 7.4 and 8.3 of
Landesman's notes \cite{Landesman2016}.
For now, only matrix factorization and substitution are needed.

For an integer \(\ell\geq1\), define

\begin{equation}\label{eq:Fell}
 F_\ell=\frac{x^\ell+y^\ell-u^\ell-v^\ell}{2}.
\end{equation}
\begin{equation}\label{eq:Segrefactor}
 F_\ell\big|_{\mathrm{Segre}}
 =\frac{(r_0^\ell-r_1^\ell)(s_0^\ell-s_1^\ell)}{2}.
\end{equation}

At \(\ell=12\), \(F_{12}=H_{\mathrm G}-H_{\mathrm E}\).
The harmonic difference becomes a product of two families of one-variable factors.
Each family has \(\ell\) linear factors; taking one factor from each family
gives a pair that can return to a linear form in the original coordinates.
The next subsection makes this pairing explicit.

Equality modulo \(Q_0\) means that the difference of two polynomials
is divisible by \(Q_0=xy-uv\). They therefore take the same values on
\(xy=uv\); this notation keeps the equal-norm constraint in view.

\Needspace{9\baselineskip}
\subsection{A two-factor example and the general pairing rule}
Take \(\ell=2\). Both lists of roots are \(1,-1\).
Pairing equal signs or opposite signs gives two product expressions:
on \(xy=uv\),
\[
 \begin{aligned}
 F_2&=\tfrac12(x+y-u-v)(x+y+u+v),\\
 F_2&=\tfrac12(x-y+u-v)(x-y-u+v).
 \end{aligned}
\]
Expansion shows that the right sides are \(F_2+(xy-uv)\) and
\(F_2-(xy-uv)\), respectively. They agree on the quadric but use
different linear factors. The following proposition extends this rule
to arbitrary \(\ell\).
Let \(\zeta_1,\ldots,\zeta_\ell\) and
\(\mu_1,\ldots,\mu_\ell\) be two lists of all \(\ell\)-th roots of unity.
\begin{proposition}\label{prop:pairings}
Over \(\C\), in the fixed coordinates \((x,y,u,v)\), all representations of
\(F_\ell\) on \(xy=uv\) as a product of \(\ell\) nonzero
homogeneous linear forms, up to factor order and individual nonzero
rescalings, are determined by bijections between the two lists of roots.
For a permutation \(\sigma\), set

\begin{equation}\label{eq:pairingforms}
 L_j=x-\mu_{\sigma(j)}u-\zeta_jv
       +\zeta_j\mu_{\sigma(j)}y,\qquad 1\leq j\leq\ell.
\end{equation}
\begin{equation}\label{eq:pairingproduct}
 F_\ell\equiv\frac12\prod_{j=1}^{\ell}L_j\pmod{xy-uv}.
\end{equation}

This gives a factorization, and every such linear-factor representation
is obtained in this way. There are \(\ell!\) equivalence classes
of factor sets; the coordinates are fixed and no further quotient
by changes of coordinates is taken.
\end{proposition}
\begin{proof}
First,

\[
 L_j\big|_{\mathrm{Segre}}
 =(r_0-\zeta_jr_1)(s_0-\mu_{\sigma(j)}s_1).
\]

Taking the product uses each root in each family once, giving
twice the expression in \eqref{eq:Segrefactor}.
To recover an identity on the quadric from an identity in the parameters,
we need the kernel of the pullback map.
The kernel of the polynomial pullback map is exactly \((xy-uv)\).
Indeed, on \(x\ne0\) choose
\(r_0=1,s_0=x,s_1=u,r_1=v/x\), obtaining \(y=uv/x\).
A polynomial with zero pullback is consequently zero in the
localized quotient ring. Since \(xy-uv\) is irreducible and does
not divide \(x\), it already divides that polynomial in the original
polynomial ring. Pullback identities thus yield identities modulo \(Q_0\).
This establishes \eqref{eq:pairingproduct}.

Conversely, suppose
\(F_\ell\equiv\kappa\prod_j\widetilde L_j\pmod{xy-uv}\),
where \(\kappa\ne0\) and each \(\widetilde L_j\) is a nonzero
homogeneous linear form. Its pullback is a nonzero polynomial
of bidegree \((1,1)\), meaning degree one in the \(r\) variables
and degree one in the \(s\) variables. In the unique factorization domain
\(\C[r_0,r_1,s_0,s_1]\), the product of these pullbacks equals
a nonzero multiple of \eqref{eq:Segrefactor}.
The latter has distinct irreducible linear factors
\(r_0-\zeta_jr_1\) and \(s_0-\mu_ks_1\).
Bidegree \((1,1)\) forces each pullback to contain exactly one
factor from each family. Every factor must occur once,
so the choices give a bijective pairing.
Pullback is injective on linear forms, since the four coordinates
map to four distinct bilinear monomials.
The original forms are therefore multiples of
\eqref{eq:pairingforms}. Different pairings give different
factor sets, yielding the count \(\ell!\).
\end{proof}

This explains why a product representation on the same quadric
need not be unique. Atkin's twelve integer-coefficient factors
give one such representation. Arbitrary pairings concern polynomial
factorization; they do not change the integral coordinates or
finite phases required by Atkin's prime formula.

For an explicit check, the pullback of \(L=Aa+Bb+Cc+Dd\) is
\((r_0,r_1)M_L(s_0,s_1)^{\mathsf T}\), where
\begin{equation}\label{eq:factormatrix}
 M_L=\frac12
 \begin{pmatrix}
 A-\iu B&C-\iu D/\sqrt3\\
 C+\iu D/\sqrt3&A+\iu B
 \end{pmatrix}.
\end{equation}
A nonzero factor separates into a linear form in \(r\) times a
linear form in \(s\) precisely when this matrix has rank one.
Substituting the coefficients of Atkin's factors checks the
specific pairing of the two lists of roots.

Each factor can also be rescaled to have real coefficients.
Writing \(\zeta=e^{\iu t}\), \(\mu=e^{\iu s}\), the corresponding \(L\) satisfies

\begin{equation}\label{eq:realpairing}
 \begin{aligned}
 \frac12 e^{-\iu(t+s)/2}L
 &=a\cos\frac{t+s}{2}+b\sin\frac{t+s}{2}\\
 &\quad{}-c\cos\frac{s-t}{2}
       +\sqrt3\,d\sin\frac{s-t}{2}.
 \end{aligned}
\end{equation}

For \(\ell=12\), \(t,s\) are integer multiples of \(\pi/6\);
the half-sum and half-difference are integer multiples of \(\pi/12\).
Choosing the pairing and rescaling its factors recovers the
coordinates of Atkin's factors. The geometric calculation
complements the symmetric-polynomial proof: it explains
pairing, while the latter verifies the integer normalization.
We now have explanations for the integer normalization and geometric pairing.
The remaining arithmetic signs are the subject of the next section.
\section{How CM determines weights and signs}
\subsection{Why the modular weight is 13}
Modular weight describes the scaling law under modular changes of variable.
Since eta has weight \(1/2\), \(\eta^{26}\) has weight 13.
On a positive-definite lattice of rank \(r\), a theta series weighted
by a homogeneous harmonic polynomial of degree \(s\) has weight
\begin{equation}\label{eq:harmonicweight}
 k=\frac r2+s.
\end{equation}
The lattice, congruence classes and finite characters determine the
precise modular group. Equation \eqref{eq:H} uses binary norms, so
\(r=2,s=12\) and \(k=1+12=13\), in agreement with \(\eta^{26}\).

The twelfth power contributes weight 12 in a binary harmonic theta series,
and the binary norm contributes weight 1. The four coordinates in
Atkin's formula compare two binary norms. A usual rank-four
theta series with harmonic degree twelve would instead have weight 14.

For an expansion in integral powers of \(q\), set
\begin{equation}\label{eq:f26}
 f_{26}(\tau)=\eta(12\tau)^{26}
            =\sum_{m\geq0}c_{26}(m)q^{12m+13}.
\end{equation}
It belongs to
\begin{equation}\label{eq:space}
 S_{13}\bigl(\Gamma_0(144),\chi_{-4}\bigr),\qquad
 \chi_{-4}(u)=\left(\frac{-4}{u}\right).
\end{equation}
Here \(S_k\) denotes cusp forms of weight \(k\),
\(\Gamma_0(144)\) is the level-144 congruence subgroup, and the
symbol is the Kronecker symbol.
Weight 13, level 144 and character \(\chi_{-4}\) place
the CM series and integer-exponent expansion in the same space.
The finite verification criterion in Section 7 will use them too.

\subsection{The exact decomposition into four Hecke theta series}
Complex multiplication, abbreviated CM, here refers to modular forms
constructed from Hecke characters of imaginary quadratic fields.
For an algebraic Hecke character \(\psi\) of infinity type 12, write
\begin{equation}\label{eq:Hecke}
 \T_\psi(\tau)
 =\sum_{\substack{\mathfrak a\subset\mathcal O_K\\
                  (\mathfrak a,\mathfrak f_\psi)=1}}
       \psi(\mathfrak a)q^{\norm(\mathfrak a)}.
\end{equation}
\Needspace{6\baselineskip}
Here \(\mathcal O_K\) is the ring of integers of \(K\),
\(\mathfrak a\) is a nonzero integral ideal, and
\(\mathfrak f_\psi\) is the conductor recording finite congruence conditions.
The condition \((\mathfrak a,\mathfrak f_\psi)=1\) restricts the sum to ideals coprime to the conductor.
Two positions in the series are worth watching:
\(\norm(\mathfrak a)\) gives the exponent of \(q\), and
\(\psi(\mathfrak a)\) gives its weight.
On principal ideals the weight contains a twelfth power of a generator
and a finite phase determined by a congruence class.
Thus norms, harmonic weights and signs appear in one expression.

Serre's decomposition can be written as
\cite{Serre1985,KrishnamoorthyDalal2024}
\begin{equation}\label{eq:CMdecomposition}
 f_{26}(\tau)=
 \frac{\T_{\psi_{1+}}(\tau)+\T_{\psi_{1-}}(\tau)
      -\T_{\psi_{2+}}(\tau)-\T_{\psi_{2-}}(\tau)}
      {32617728}.
\end{equation}
Here \(K_1=\Q(\sqrt{-3})\) and \(K_2=\Q(\iu)\).
The conductors are \(4\sqrt{-3}\,\mathcal O_{K_1}\) and
\(6\mathcal O_{K_2}\), of norms 48 and 36 respectively.
The associated modular levels are therefore both 144:
\(3\cdot48=4\cdot36=144\). The denominator is \(32C\).

The finite phases use the following conventions.
In this subsection \(x,y,u,v\) are generator and residue-class parameters,
used separately from the quadric coordinates in Section 3.
For an ideal in \(K_1\) coprime to the conductor, choose a generator
\(\alpha=x+y\sqrt{-3}\) with \(x,y\in\Z\),
\(x+y\equiv1\pmod2\), and \(x\equiv1\pmod3\). Set
\begin{equation}\label{eq:psi1}
 \psi_{1\pm}((\alpha))=(-1)^{(x\mp y-1)/2}\alpha^{12}.
\end{equation}
\Needspace{9\baselineskip}
For a generator \(\beta\) in \(K_2\), choose
\(u\in\Z/2\Z,\ v\in\Z/8\Z\) such that
\[
 \beta\equiv\iu^u\pmod{2\mathcal O_{K_2}},\qquad
 \beta\equiv(1-\iu)^v\pmod{3\mathcal O_{K_2}},
\]
and set
\begin{equation}\label{eq:psi2}
 \psi_{2\pm}((\beta))=(-1)^{3u}(\pm\iu)^{3v}\beta^{12}.
\end{equation}
At a prime splitting in both fields, pairing an ideal with its conjugate
produces \(\Repart(\alpha^{12})\) and \(\Repart(\beta^{12})\).
This agrees with the algebraic reconstruction in
Proposition \ref{prop:harmonic}. Binary quadratic-form expressions
also appear in \cite{ChanCooperToh2006}.
Since \(f_{26}\) is a linear combination of four eigenforms, its
coefficients must not be treated without proof as the multiplicative
coefficients of a single eigenform.

\Needspace{10\baselineskip}
\subsection{How finite phases determine the common sign}
Write \([q^p]F\) for the coefficient of \(q^p\) in \(F\).
The proof has two steps: read the signs from the characters of the two fields,
then use coefficient integrality to show that these signs are equal.
\begin{proposition}\label{prop:CMAtkin}
Equation \eqref{eq:CMdecomposition}, together with coefficient
integrality, implies \eqref{eq:AtkinH}, including the common sign
\((-1)^a\).
\end{proposition}
\begin{proof}
Take a prime as in \eqref{eq:representations}.
Since \(p\equiv1\pmod4\), \(c\) is odd and \(d\) is even.
Change the sign of \(c\) to obtain \(x\equiv1\pmod3\), and put
\[
 s_{\mathrm E}=(-1)^{(x+d-1)/2}.
\]
Both characters in \eqref{eq:psi1} have this same finite sign on the
two conjugate prime ideals. Hence
\([q^p]\T_{\psi_{1\pm}}=2s_{\mathrm E}H_{\mathrm E}\).

In the Gaussian field, take a generator with odd real part and even
imaginary part. If \(a\) is even, the real part comes from \(b\) and the
imaginary part from \(a\). Since \(3\mid a\), the generator is
\(\pm1\) modulo 3, so \(v=0,4\) and the phase in \eqref{eq:psi2}
is \(+1\). If \(a\) is odd, the generator is \(\pm\iu\) modulo 3,
so \(v=2,6\) and the phase is \(-1\).
In either case \(u=0\) modulo 2. Thus
\(s_{\mathrm G}=(-1)^a\) and
\([q^p]\T_{\psi_{2\pm}}=2s_{\mathrm G}H_{\mathrm G}\).
Exchanging the two square coordinates preserves the twelfth-power
real part, since \(\iu^{12}=1\).

Substitution into \eqref{eq:CMdecomposition} gives
\begin{equation}\label{eq:CMprimephase}
 \tau_{26}(p)=
 \frac{s_{\mathrm E}H_{\mathrm E}-s_{\mathrm G}H_{\mathrm G}}{8C}.
\end{equation}
Moreover, \(H_{\mathrm E}\equiv H_{\mathrm G}\equiv1\pmod3\),
using \(3\nmid c\) in the first case and \(3\mid a,\ 3\nmid b\)
in the second. The left side is integral and \(3\mid C\), so
\(s_{\mathrm E}\equiv s_{\mathrm G}\pmod3\).
As both signs belong to \(\{1,-1\}\), they are equal and both are
\((-1)^a\). This yields \eqref{eq:AtkinH}; applying
Proposition \ref{prop:harmonic} then recovers the twelve-factor formula.
\end{proof}
The twelfth-power real parts thus come from generators, the signs from
finite congruence phases, and equality of the two signs from the integrality
supplied by the Euler product.


\section{From prime coefficients to all coefficients}
\subsection{The full Chan--Cooper--Toh coefficient formula}
Prime indices have canonical unique representations, whereas composite
indices may have several. For all coefficients one should use a full
sum rather than substitute an arbitrary quadruple into Atkin's product.
Define

\begin{equation}\label{eq:Phi}
 \Phi(X,Y)=X^6-66X^5Y+495X^4Y^2-924X^3Y^3
             +495X^2Y^4-66XY^5+Y^6.
\end{equation}
\begin{equation}\label{eq:Phiharmonic}
 \Phi(s^2,t^2)=\Repart(s+\iu t)^{12}.
\end{equation}

Comparing coefficients in Theorem 3 of Chan--Cooper--Toh yields their
Corollary 4: for every integer \(m\geq0\), \cite{ChanCooperToh2006}
\begin{align}
 16C\,c_{26}(m)
 &=\sum_{\substack{\alpha^2+\beta^2=24m+26\\
                  \alpha,\beta\equiv1\ (6)}}
    (-1)^{(\alpha+\beta-2)/6}
       \Phi\left(\frac{\alpha^2}{2},\frac{\beta^2}{2}\right)
       \notag\\
 &\quad{}+
   \sum_{\substack{\gamma^2+3\delta^2=36m+39\\
                  \gamma\equiv0\ (6),\ \delta\equiv1\ (6)}}
    (-1)^{(\gamma+\delta-1)/6}
       \Phi\left(\frac{\gamma^2}{3},\delta^2\right).
 \label{eq:CCTcoefficients}
\end{align}

Read \eqref{eq:CCTcoefficients} as a procedure: enumerate integer pairs
satisfying the norm and congruence conditions, assign the signs and harmonic weights,
then add and divide by \(16C\).
Here \(\alpha,\beta,\gamma,\delta\) are integer summation variables.
Both sums include negative coordinates satisfying the congruences. Restricting \(\alpha,\beta\) to positive
integers, for example, would change the formula.
The constant \(16C=16308864\) differs from the \(32C\) of the CM
decomposition because the coordinate scaling and representation counts differ.

The two norms are \(2(12m+13)\) and \(3(12m+13)\).
The harmonic weights have the integer-arithmetic expressions
\[
 \Phi\left(\frac{\alpha^2}{2},\frac{\beta^2}{2}\right)
 =\frac{\Repart(\alpha+\iu\beta)^{12}}{64},\qquad
 \Phi\left(\frac{\gamma^2}{3},\delta^2\right)
 =\frac{\Repart(\gamma+\iu\sqrt3\,\delta)^{12}}{729}.
\]
The full double series and the CM decomposition describe the same
coefficient, although their integer coordinates must not be mixed.
For prime \(12m+13\), canonical representations simplify the sums.
For composite indices, all representations and finite signs are
needed to capture cancellation.

Section 7 computes the same coefficients using the Euler-product recurrence
and this sum, allowing the indexing and normalization to be checked directly.
\subsection{CM support and lacunarity}
A power series is lacunary when the nonzero coefficients have density zero:
\begin{equation}\label{eq:lacunary}
 \lim_{M\to\infty}
 \frac{\#\{0\leq m\leq M:c_{26}(m)\ne0\}}{M+1}=0.
\end{equation}
This definition concerns the unscaled series \eqref{eq:Euler26},
not merely the empty positions imposed by \(n\equiv1\pmod{12}\)
in \eqref{eq:f26}.

Serre proved that the positive even powers of eta are lacunary precisely
for exponents
\begin{equation}\label{eq:Serrelist}
 2,\;4,\;6,\;8,\;10,\;14,\;26 .
\end{equation}
This classification is due to Serre \cite{Serre1985}. Lacunarity allows
infinitely many nonzero terms; their proportion among all indices must tend to zero.
An inert prime with odd exponent obstructs an ideal-norm representation
in the relevant quadratic field. The two CM fields impose their
respective support restrictions, and their linear combination may
produce additional cancellation. The full CM decomposition records both support restrictions
and possible cancellation.

\subsection{How coefficients at composite indices vanish}
Let \(n=12m+13\), and let \(v_\ell(n)\) denote the exponent of a prime
\(\ell\) in \(n\). Either of the following conditions implies
\(c_{26}(m)=0\) \cite{Serre1985}:
\begin{enumerate}
 \item There is a prime \(\ell\equiv3\pmod4\) with \(v_\ell(n)\) odd,
 and a prime \(\ell'\equiv2\pmod3\) with \(v_{\ell'}(n)\) odd.
 The two primes may coincide.
 \item The integer \(n\) is a square, and all its prime factors satisfy
 \(\ell\equiv11\pmod{12}\).
\end{enumerate}
The first condition blocks both groups of CM coefficients through
inert primes, making all four terms of \eqref{eq:CMdecomposition}
vanish. The second is different: norm representations remain possible,
but the weighted terms cancel.

Indeed, for the coefficient function \(A(n)\) of any of the four
normalized eigenforms, and \(\ell\nmid144\),
\begin{equation}\label{eq:Heckerecurrence}
 A(\ell^j)=A(\ell)A(\ell^{j-1})
 -\chi_{-4}(\ell)\ell^{12}A(\ell^{j-2}),\qquad A(1)=1.
\end{equation}
For \(\ell\equiv11\pmod{12}\), both fields are inert, so
\(A(\ell)=0\) for every form, and \(\chi_{-4}(\ell)=-1\).
It follows that
\[
 A(\ell^{2s})=\ell^{12s},\qquad A(\ell^{2s+1})=0.
\]
For a square \(n\) of the second type, all four terms equal \(n^6\);
two positive and two negative terms cancel. The first example is
\begin{equation}\label{eq:firstzero}
 n=11^2=121,\qquad m=9,\qquad c_{26}(9)=0.
\end{equation}
There are still norm representations, for example
\[
 121=11^2+0^2=11^2+3\cdot0^2.
\]
These show that both fields admit norm representations; the actual coefficient
sums must also satisfy their respective congruence conditions.
The zero coefficient results from cancellation of the four weighted terms.
This is consistent with prime nonvanishing in Corollary \ref{cor:prime}.
The existence of norm representations and the nonvanishing of a
signed weighted coefficient are different questions.


\subsection{Further reading: converses to the vanishing criteria}
Krishnamoorthy and Dalal studied converses to these vanishing
criteria in 2024. Their Theorems 1.4 and 1.5 concern
\(c_{26}(25u+1)\) and \(c_{26}(49u+3)\), respectively.
For primes \(p\equiv1\pmod{12}\) dividing \(12u+1\), they impose
exponent restrictions excluding \(4\pmod5\) and \(6\pmod7\),
respectively \cite{KrishnamoorthyDalal2024}.
Under these hypotheses the two inert-prime obstructions give
the corresponding necessary and sufficient criterion.
These are conditional partial converses, rather than an unrestricted
classification of vanishing at all composite indices.

\section{The connection with Jacobi specialization}
A Jacobi form has an elliptic variable \(z\) in addition to \(\tau\).
Differentiating in \(z\) and setting \(z=0\) can return an expression
depending only on \(\tau\).
Wagner constructs the corresponding \(\phi_{26}(z;\tau)\) from four CM theta terms.
Theorem 1.3 gives, in its normalization of the elliptic variable,
\cite{Wagner2023}

\begin{equation}\label{eq:Wagner}
 \left.D_z^{12}\phi_{26}(z;\tau)\right|_{z=0}
 =\frac{\eta(\tau)^{26}}{2^8\cdot3^4\cdot11^2\cdot13},
 \qquad D_z=\frac{1}{2\pi\iu}\frac{\partial}{\partial z}.
\end{equation}

The expression has a binary-quadratic-form sum representation
and a representation as a sum of infinite products.
The subsequent remark distinguishes this from the desired single
theta product, whose expected shape is
\begin{equation}\label{eq:expectedblock}
 \frac{\prod_{j=1}^{12}\vartheta_1(L_j(z);\tau)}
      {\eta(\tau)^{10}}.
\end{equation}

Here \(\vartheta_1\) is the odd Jacobi theta function and the \(L_j\) are linear parameters.
Consider their variation along a chosen direction. In the usual normalization,
\(\vartheta_1'(0;\tau)=2\pi\eta(\tau)^3\).
If each \(L_j\) is nonzero on the chosen direction,
\eqref{eq:expectedblock} has weight \(12/2-10/2=1\),
vanishing order twelve, and a leading term containing
\(\eta^{36-10}=\eta^{26}\).
This explains the local origin of the twelve theta factors and
the tenth power in the denominator.
The local expansion does not establish a global Jacobi
transformation law, lattice index or holomorphy.
Nor does the polynomial pairing in Section 3 by itself produce
such a Jacobi expression.

The connection involves both local expansions and global transformation laws.
The local expansion explains the twelve factors and the appearance of \(\eta^{26}\).
A global identity also requires a lattice index, Jacobi transformations
and a global zero structure.
The polynomial pairing in Section 3 gives an intuitive guide to the local structure;
these analytic conditions must also be checked to obtain a product identity for functions.

\section{How to check the formulas yourself}
\subsection{Cross-checking three coefficient computations}
Begin with the Euler product, which gives a direct computational method.
Taking the logarithmic derivative of \eqref{eq:Euler26} yields
\[
 q\frac{\diff}{\diff q}\log E_{26}(q)
 =-26\sum_{j\geq1}\sigma_1(j)q^j.
\]
Multiplying by \(E_{26}(q)\) and comparing coefficients gives the recurrence
\begin{equation}\label{eq:coefficientrecursion}
 c_{26}(0)=1,\qquad
 m c_{26}(m)=-26\sum_{j=1}^m\sigma_1(j)c_{26}(m-j),
 \qquad \sigma_1(j)=\sum_{t\mid j}t.
\end{equation}
Here \(\sigma_1(j)\) is the sum of the positive divisors of \(j\).
For \(m=1\), it gives \(c_{26}(1)=-26\).
For \(m=2\), it gives \(2c_{26}(2)=-26(-26+3)=598\), so the coefficient is 299.
Starting with \(c_{26}(0)=1\) and proceeding in order computes every coefficient needed.

The accompanying \texttt{verify\_d26.py} compares this recurrence,
the full Chan--Cooper--Toh sum and Atkin's prime product using integer arithmetic.
For \(0\leq m\leq1000\), all 1001 double-series coefficients agree with the recurrence;
the 345 cases with prime \(12m+13\) also agree with Atkin's formula.
The script checks finite phases and the first vanishing coefficient as well.
Table \ref{tab:samples} gives several points for checking by hand or by script.

The symbolic option checks divisibility of \(H_{\mathrm G}-H_{\mathrm E}+8P\)
by \(Q_0\), the two harmonicity identities and the Segre pullback.
All twelve Atkin factor matrices have rank one, with each root in each list
occurring once, in agreement with Proposition \ref{prop:pairings}.
These computations help locate errors in constants, signs and indexing;
the general statements are proved by the preceding algebraic arguments.
\begin{table}[H]
\centering
\caption{Exact checks at selected prime indices}\label{tab:samples}
\begin{tabular}{rrrrrrr}
\toprule
\(p\) & \(m=(p-13)/12\) & \(a\) & \(b\) & \(c\) & \(d\) & \(c_{26}(m)\)\\
\midrule
13  & 0  & 3 & 2  & 1  & 2 & 1\\
37  & 2  & 6 & 1  & 5  & 2 & 299\\
61  & 4  & 6 & 5  & 7  & 2 & 7475\\
73  & 5  & 3 & 8  & 5  & 4 & \(-13754\)\\
97  & 7  & 9 & 4  & 7  & 4 & 132756\\
109 & 8  & 3 & 10 & 1  & 6 & \(-276575\)\\
157 & 12 & 6 & 11 & 7  & 6 & \(-519961\)\\
181 & 14 & 9 & 10 & 13 & 2 & 829725\\
193 & 15 & 12 & 7 & 1  & 8 & \(-2215486\)\\
\bottomrule
\end{tabular}
\end{table}

\subsection{A finite and rigorous modular-form verification criterion}
If a candidate is already proved to lie in the same holomorphic
modular-form space \eqref{eq:space}, a Sturm-type bound reduces
identity checking to finitely many coefficients.
The group index is
\begin{equation}\label{eq:indexgroup}
 [\mathrm{SL}_2(\Z):\Gamma_0(144)]
 =144\left(1+\frac12\right)\left(1+\frac13\right)=288.
\end{equation}
For weight 13, the bound is
\begin{equation}\label{eq:Sturm}
 B=\left\lfloor\frac{13}{12}\cdot288\right\rfloor=312.
\end{equation}
Two forms in the same space are equal if their coefficients of
\(q^n\) agree for every \(0\leq n\leq312\).
For the target support, the relevant nonzero positions are
\(n=13,25,\ldots,301\), namely \(c_{26}(0),\ldots,c_{26}(24)\);
the other positions in this range must also be checked to be zero.

The bound can be understood without special software.
For the difference \(h\), the equality \(\chi_{-4}^2=1\) makes
\(h^2\) a holomorphic weight-26 modular form with trivial character.
If the first possible nonzero term of \(h\) is \(q^{313}\) or later,
the vanishing order of \(h^2\) at infinity is at least 626, exceeding
the valence bound \(26\cdot288/12=624\). Hence \(h\) vanishes identically.
The hypothesis of membership in the same holomorphic modular space
is essential. Without transformations and holomorphy at the other
cusps, extensive numerical agreement cannot invoke this bound.
Checking only prime coefficients is not this criterion.


\subsection{Connecting the different expressions}
The formulas can now be read along one continuous thread.
Recognize Atkin's twelve factors as an equal-norm harmonic difference,
use the parametrization to see the two lists of linear factors and their pairing,
use CM characters to supply finite signs, and finally use the full sums
and Hecke recurrence to understand multiple representations and cancellation.
Each section describes a different layer of the same coefficient problem.

Checks can follow the same order: test the constant and sign at \(p=13\),
compute low coefficients by recurrence, verify the polynomial identity on the quadric,
and then use a finite criterion once the modular-form space has been established.
This keeps the relationships between the formulas visible while making clear
what each step verifies.

\appendix
\section{Computation and compilation}
Both UTF-8 source files embed their references and need no external figures
or BibTeX database. Compile twice with XeLaTeX to update cross-references.
The Chinese source \texttt{dyson\_d26.tex} uses \texttt{ctex} when available;
otherwise it tries Songti SC, Noto Serif CJK SC or FandolSong.
The English source is \texttt{dyson\_d26\_en.tex}.

\Needspace{10\baselineskip}
Both versions share \texttt{verify\_d26.py}. Run
\begin{lstlisting}
python3 verify_d26.py --limit 1000
# Optional: install SymPy for symbolic checks.
python3 verify_d26.py --limit 1000 --symbolic
\end{lstlisting}
The first command needs only the Python standard library.
The second needs SymPy for polynomial remainders, harmonicity,
Segre substitution and factor matrices. These steps use exact arithmetic.
For coefficients through \(q^M\), the Euler product can be truncated to \(j\leq M\).

\Needspace{19\baselineskip}
\begingroup
\small
\linespread{1.02}\selectfont
\begin{thebibliography}{9}
\setlength{\itemsep}{.35em}
\setlength{\parskip}{0pt}
\raggedright
\bibitem{Dyson1972}
F.~J.~Dyson.
\textit{Missed opportunities}.
Bulletin of the American Mathematical Society \textbf{78}(5) (1972),
635--652. The prime formula discussed here is in Note 14.

\href{https://doi.org/10.1090/S0002-9904-1972-12971-9}{doi:10.1090/S0002-9904-1972-12971-9}.
\href{https://wucj.lab.westlake.edu.cn/Others/Dyson_missedopportunity.pdf}{Public scan of the original article}.
\bibitem{Serre1985}
J.-P.~Serre.
\textit{Sur la lacunarité des puissances de \(\eta\)}.
Glasgow Mathematical Journal \textbf{27} (1985), 203--221.
\href{https://doi.org/10.1017/S0017089500006194}{doi:10.1017/S0017089500006194}.

\bibitem{ChanCooperToh2006}
H.~H.~Chan, S.~Cooper and P.~C.~Toh.
\textit{The 26th power of Dedekind's \(\eta\)-function}.
Advances in Mathematics \textbf{207}(2) (2006), 532--543.
\href{https://math.nie.edu.sg/pctoh/archive/26thpowerofeta.pdf}{Author manuscript}.

\bibitem{Wagner2023}
I.~Wagner.
\textit{Jacobi forms with CM and applications}.
Journal of Number Theory \textbf{246} (2023), 15--48.
\href{https://doi.org/10.1016/j.jnt.2022.11.014}{doi:10.1016/j.jnt.2022.11.014}.
\href{https://arxiv.org/abs/2208.01807}{arXiv:2208.01807}.

\bibitem{KrishnamoorthyDalal2024}
S.~Krishnamoorthy and T.~Dalal.
\textit{On the vanishing of coefficients of \(\eta^{26}\)}.
Preprint (2024).
\href{https://arxiv.org/abs/2411.14990}{arXiv:2411.14990}.


\bibitem{Landesman2016}
A.~Landesman.
\textit{Math 137 Notes: Undergraduate Algebraic Geometry}.
Notes from J.~Harris's course, Harvard University, Spring 2016,
Sections 7.4 and 8.3.
\href{https://people.math.harvard.edu/~landesman/assets/harris-undergrad-algebraic-geometry-notes.pdf}
{Public course notes}.
\end{thebibliography}
\endgroup
\end{document}
